\answer{ex:19:1}{(a)~Une synapse donne $6.7\mV$ ($R=66.7$~M$\Omega$): pour atteindre le seuil tout court, il en faut au moins $4$ (trois ne suffisent qu'asymptotiquement); $8$ en $10\ms$; $14$ en $5\ms$. (b)~$0.1\ms\ll\tau_m$, correction due à la fuite $\sim0.25\%$.}
\answer{ex:19:2}{(a)~$1.75\cdot10^{10}$ (un quart des mélangeurs répond à tout), $4.4\cdot10^9$ et $0.06$ (pour $k=40$, presque jamais de réponse). (b)~D'un facteur $10^3$: $2\cdot10^5$ contre $200$.}
\answer{ex:19:3}{(a)~$C(2n,n)=2^{2n-1}$ par symétrie des coefficients binomiaux. (b)~$0.93$ et $0.09$; largeur de la transition $\sim\sqrt{2n}$ en $P$, relative $\sim1/\sqrt n$.}
\answer{ex:19:4}{Avec une seule dendrite, $V_{\text{corps}}<\kappa E_E<\theta$ quel que soit le nombre d'entrées; pour $g_0=0.5g_L$, il faut $n\ge3$ sur chaque dendrite.}
\answer{ex:19:5}{$I\ge C\sigma_V/\sigma_t=15$~nA, soit $150$ synapses synchrones, ce qui est irréaliste; au petit neurone, $1$~nA suffirait, et avec une synapse de $4$~nA la gigue vaut $5$~µs.}
\answer{ex:19:6}{Maximum à l'extrémité proche: $0.03\,\tau_m$ et $0.97$ ($L=0.2$); $0.32\,\tau_m$ et $0.66$ ($L=1$); $0.79\,\tau_m$ et $0.33$ ($L=2$). L'estimation~\eqref{eq:dendriteload} par la capacité totale n'est valable que pour $L\ll1$.}
\answer{ex:19:7}{Question ouverte. À $m$ fixé, le temps de réponse croît linéairement avec le nombre $P$ d'associations mémorisées; à fraction fixée $m=\varphi n$, il ne dépend plus de $n$, et la lenteur d'un grand neurone découle de la sommation de nombreuses entrées, et non de la taille en tant que telle.}
